% [AUTHOR REWRITE] fit claim
The extension also changes how well the model predicts the apartment premium, though less than the aggregate suggests (Figure~\ref{fig:att_fit}). It removes the level bias in the predicted premia, the mean residual falling from \attResidCoreApt{} to \attResidExtApt{} log points, but that zero holds by construction, because $\hat\psi_\delta$ is estimated from the same ratio. The root mean squared error falls from \attRmseCoreApt{} to \attRmseExtApt{}, yet net of the level the residual dispersion is unchanged (\attSdResidCoreApt{} against \attSdResidExtApt{}) and the correlation between predicted and observed premia moves only from \attCorrCoreApt{} to \attCorrExtApt{}. The extension therefore explains why silent apartments sell above verified efficient ones, not which markets carry the largest premia. In houses the two models coincide. The evidence for the extension rests on the ceilings of Figures~\ref{fig:att_belief} and~\ref{fig:att_ceilings} (the latter in Appendix~\ref{app:attention}) and the contrast of Figure~\ref{fig:att_contrast}, not on the aggregate fit.
